\documentclass[11pt]{article}
\usepackage[letterpaper,margin=1in]{geometry}
\usepackage{amsmath,amssymb,amsthm,booktabs,longtable,array}
\usepackage[colorlinks=true,linkcolor=black,urlcolor=black,citecolor=black]{hyperref}
\usepackage{url}
\newtheorem{theorem}{Theorem}[section]
\newtheorem{lemma}[theorem]{Lemma}
\newtheorem{proposition}[theorem]{Proposition}
\newtheorem{corollary}[theorem]{Corollary}
\theoremstyle{definition}\newtheorem{definition}[theorem]{Definition}
\newcommand{\E}{\mathbb E}
\newcommand{\Prob}{\mathbb P}
\newcommand{\HH}{\mathcal H}
\newcommand{\XX}{\mathcal X}
\newcommand{\KL}{\operatorname{KL}}
\newcommand{\JS}{\operatorname{JS}}
\newcommand{\ind}{\mathbf 1}
\newcommand{\err}{\widehat L}
\newcommand{\decl}[1]{{\small\nolinkurl{#1}}}
\setlength{\emergencystretch}{2em}
\setcounter{tocdepth}{1}
\title{Conditional mutual information and empirical error in proper VC learning}
\author{Solved Conjectures Project}
\date{November 2026}
\hypersetup{pdftitle={Conditional mutual information and empirical error in proper VC learning},
 pdfauthor={Solved Conjectures Project},pdfsubject={Conditional mutual information, proper learning and measurable selection}}
\begin{document}
\maketitle
\begin{abstract}
For a nonempty binary hypothesis class of VC dimension at most $d$ and
$n\ge d\ge1$, we construct a proper randomized learner whose expected empirical error on every sample is within
$583\sqrt{d/n}$ of the class optimum. On every fixed pool of $2n$ labeled
observations, its selector information, averaged over uniform orderings of the
pool, is at most $19{,}712d$ nats. This bound concerns the identity of the
returned classifier. Exchangeability gives the iid bound whenever the
pool statistic is almost everywhere measurable. Measurability of trace realizability, representative extensions and output
events gives a single Markov learner for every iid law.
This includes arbitrary classes on countable discrete domains and classes
with a countable measurable subclass realizing all their finite traces. We identify the extra hypothesis
relative to Steinke and Zakynthinou's Conjecture~6 and exhibit a VC-one
obstruction to an unrestricted Borel-kernel interpretation.
\end{abstract}

\section{Introduction}\label{sec:statement}
Let $\HH$ be a nonempty class of functions from an instance set $\XX$ to
$\{0,1\}$. Its VC dimension is the largest number of distinct points on
which the class realizes every binary labeling. For a sample
$z=((x_i,y_i))_{i=1}^n$, set
\[
 \err_z(h)=\frac1n\sum_{i=1}^n\ind\{h(x_i)\ne y_i\}.
\]
A proper learner $A$ returns a member of $\HH$. Its internal randomness
is independent of the data. Let $Z=(Z_{i,b})_{i\le n,b\in\{0,1\}}$ be a pool
of $2n$ independent observations with common law $D$, and let $S$ be an
independent vector of $n$ fair bits. Write $Z_S=(Z_{i,S_i})_{i=1}^n$.
Steinke and Zakynthinou's original conditional mutual information is
$I(S;A(Z_S)\mid Z)$ \cite[Definition 3]{problem}. Even two classifiers that
agree on every training point can be different outputs in this experiment.

Steinke and Zakynthinou \cite[Conjecture 6]{problem} ask for an absolute $c$ such that for every binary class
of VC dimension $d\ge1$ and every $n\ge d$, there exists one proper $A$ with
\begin{align}
  &\forall D,\quad I(S;A(Z_S)\mid Z)\le cd, \qquad
    \forall z,\quad \E_A\err_z(A(z))
       \le \inf_{h\in\HH}\err_z(h)+c\sqrt{d/n}.
       \label{eq:conj}
\end{align}
The learner is chosen before $D$. The authors' preceding theorem gives an
empirical risk minimizer with a global tie-breaking rule with information at most
$d\log n+2$ \cite[Theorem 5]{problem}. Here we allow a small empirical excess
to obtain a bound independent of $n$.

\paragraph{Fixed-pool information.}
For a deterministic labeled pool $t$, let $Q_s$ be the law of $A(t_s)$
and $\bar Q=2^{-n}\sum_s Q_s$. Define
\begin{equation}
  F_A(t)=2^{-n}\sum_s\KL(Q_s\,\|\,\bar Q).
  \label{eq:fixed}
\end{equation}
All these laws have finite support for our learner. Their atoms are actual
hypotheses in $\HH$. Writing $Y=A(t_S)$, the joint mass of $(S,Y)$ at $(s,h)$ is
$2^{-n}Q_s(h)$ and its marginal product has mass $2^{-n}\bar Q(h)$.
Their log ratio is the one in \eqref{eq:fixed}, so this formula is precisely
the selector information on that pool. We use natural logarithms and the
usual zero-mass convention.

\begin{theorem}[Fixed pools]\label{thm:pool}
Fix $1\le d\le n$ and a nonempty binary class of VC dimension at most $d$.
There is a proper randomized learner $A$ such that
\begin{align}
 \E_A\err_z(A(z))&\le\inf_{h\in\HH}\err_z(h)+583\sqrt{d/n}
                      &&\text{for every labeled sample }z, \label{eq:accuracy}\\
 \E_{\pi}\,F_A(t\circ\pi)&\le19{,}712d
                      &&\text{for every labeled pool }t. \label{eq:pool}
\end{align}
Here $\pi$ uniformly permutes the $2n$ positions. The order is fixed when
computing $F_A$ and averaged afterwards. Repeated instances and conflicting
labels are allowed.
\end{theorem}

\begin{theorem}[Population form]\label{thm:iid}
For the same $A$, under every iid law for which $F_A$ is almost everywhere
measurable, $F_A$ is integrable and $\E_Z F_A(Z)\le19{,}712d$.
If the chosen representatives satisfy Definition~\ref{def:selection}, the
output law is a measurable Markov kernel, and this inequality holds for
every iid law. Under this kernel hypothesis, if output singletons are measurable, $\E_Z F_A(Z)$ equals
the original conditional mutual information defined using measure-theoretic
KL. With a coarser output sigma-algebra its observable KL is bounded above by
the complete-name information on each pool.
\end{theorem}

Here a Markov kernel means a rule assigning each sample a probability law on
$\HH$, with the probability of each measurable output event a measurable
function of the sample. The finite-pool statement uses no measurable-space
structure. Conjecture~6 does not specify structures on all possible classes.
Theorem~\ref{thm:iid} proves \eqref{eq:conj} under an explicit selection
hypothesis. Section~\ref{sec:scope} explains why that qualification matters.

\paragraph{Related work.}
Steinke and Zakynthinou also ask realizable versions with sharper empirical
requirements. Haghifam, Dziugaite, Moran and Roy give negative results for
those proper-learning variants \cite{haghifam}; their required empirical
error on realizable samples is smaller than the allowance in
\eqref{eq:conj}. Hanneke and Wang obtain linear \emph{evaluated} information
for an improper majority learner in the realizable case \cite{hanneke}.
Evaluated information measures the loss vector on the supersample. Our
channel keeps the complete hypothesis name, including information invisible
in that vector.

\paragraph{Proof outline.}
A learner can randomize over the distinct prediction vectors on its training
sample. To be proper it must also choose a named hypothesis for each vector.
We use the first such hypothesis in one global order. For accuracy we must
keep the randomization close to an empirical minimizer. For information we
must control how much the entire output law changes on deleting an example.
The same objective serves both purposes: it rewards low weighted loss and
high entropy at several sizes of feature subsets. The proof below first
compares its value with a point mass, then compares its minimizers before
and after deletion.

\section{Construction of the learner}\label{sec:learner}
Fix a well-order $\prec$ of $\HH$ once for all datasets. For an ordered feature
tuple $x=(x_i)_{i\in I}$, let $V(x)=\{(h(x_i))_{i\in I}:h\in\HH\}$.
A vector $v\in V(x)$ is a \emph{trace}; write $g_x(v)$ for its first
realizing hypothesis. These names are coherent under restriction. Indeed,
if $g$ is first among the hypotheses realizing a trace on a subtuple, any
earlier hypothesis with the same \emph{full} trace as $g$ would also realize
that subtuple trace. Thus $g$ is the first hypothesis for its full trace too.
This fact puts the outputs of a sample and all its deletions on one common
alphabet of actual hypotheses. For example, suppose $\HH$ consists of the
zero function and the indicator of a point $b$, in that order. On two
distinct features $a,b$, their traces are $(0,0)$ and $(0,1)$. Delete $b$
and both traces become $(0)$; its chosen name is the zero function. A proof
that compares only surviving predictions would miss the probability that
the original learner returned the singleton classifier. The deletion
argument keeps track of precisely this mass.

For $t=n/d$, let $J$ be the unique integer with $16^J\le t<16^{J+1}$.
Set
\[
 \rho=21/20,\qquad m_j=d4^j,\qquad b_j=\sqrt{t}\,\rho^{-j},\qquad
 a_j=b_j-b_{j+1}\ (0\le j\le J),\qquad \tau=b_{J+1}.
\]
In particular $m_J^2\le nd$ and $m_J\le n$. The following definition also
applies to a working block of $n+1$ labeled examples, always using these
parameters for the original $n$.

\begin{enumerate}
\item On the given sample draw independent $U_i\sim\operatorname{Unif}[0,2]$.
For $v\in V(x)$ put $L_U(v)=\sum_i U_i\ind\{v_i\ne y_i\}$.
\item For a law $q$ on $V(x)$, denote its marginal trace on $M\subseteq I$
by $q_M$ and its entropy by $H(q)=-\sum_v q(v)\log q(v)$.
Choose the law minimizing
\begin{equation}
 \Phi_U(q)=\sum_v q(v)L_U(v)-\tau H(q)
             -\sum_{j=0}^J a_j\,\E_{|M|=m_j} H(q_M).
 \label{eq:objective}
\end{equation}
Each $M$ is uniform among subsets of that size; the average is part of the
objective.
\item Draw a trace $v$ from this law and output $g_x(v)$.
\end{enumerate}

The simplex is finite-dimensional and compact. The objective is continuous
and strictly convex, since $\tau>0$, negative entropy is strictly convex,
and each restricted negative entropy is convex. The unique minimizer gives
positive mass to every trace: if a trace has zero mass, mixing in a point
mass of size $\varepsilon$ there increases full entropy by
$\varepsilon\log(1/\varepsilon)+O(\varepsilon)$. Concavity bounds the possible
loss of each restricted entropy by $\varepsilon H(q_M)$, and the linear cost
changes by $O(\varepsilon)$. Hence mixing would improve the objective.
The entropy terms are therefore differentiable at the minimizer. For fixed features and labels the minimizing law is continuous in
the weights: every convergent subsequence of minimizers over converging
weights minimizes the limiting objective, which has a unique minimizer.

\section{Empirical error}\label{sec:accuracy}
Fix a sample and abbreviate $\ell_i(v)=\ind\{v_i\ne y_i\}$ and
$L(v)=\sum_i\ell_i(v)$. Let $q$ minimize \eqref{eq:objective} and let $v_*$ be
a trace minimizing $L$. The point mass at $v_*$ has zero entropy at every
scale. Using it as a competitor and undoing the weights gives
\begin{align}
 \E_{v\sim q}L(v)-L(v_*)&\le
  \tau\log|V(x)|+\sum_j a_j\E_{|M|=m_j}\log|V_M|
      +2\max_{v\in V(x)}\left|\sum_i(U_i-1)\ell_i(v)\right| .
 \label{eq:accbudget}
\end{align}
Here $V_M=\{v|_M:v\in V(x)\}$. The factor of two accounts for the change of weights on both $q$
and $v_*$. Let $R=\sqrt{nd}$. The three terms on the right, in expectation,
are bounded by $320R$, $81R$ and $2\cdot91R$ respectively.

\paragraph{The deterministic terms.}
If a set of positions is shattered, their features must be distinct:
two equal features cannot receive opposite predictions. Their feature values
are then shattered by $\HH$. Thus the VC dimension bound applies to trace
families even when the training tuple contains repetitions.
On $m\ge d$ positions the Sauer--Shelah bound \cite{sauer} gives
$\log |V_M|\le d(1+\log(m/d))$ (a proof is in Appendix~\ref{app:emp}).
Since $t<16^{J+1}$, with $k=J+1$ we have
$1+\log t\le1+16k\le320\rho^k$ by Bernoulli's inequality. Multiplying
by $d\tau$ gives $320R$. At scale $m_j$, $1+\log(m_j/d)\le1+4j$.
The identity
\[
 (b_j-b_{j+1})(1+4j)=b_j(81+4j)-b_{j+1}(81+4(j+1))
\]
follows from $b_j/b_{j+1}=21/20$. Summing gives the second bound.

\paragraph{The fluctuation term.}
The centered weights have the law of independent fair signs times independent
amplitudes in $[0,1]$. Conditional on the amplitudes, the loss differences
between two traces depend only on positions where they disagree. A packing
argument bounds the number of distinct traces needed at each Hamming scale.
Summing the expected sign fluctuations across those scales gives
\[
 \E_U\max_v\left|\sum_i(U_i-1)\ell_i(v)\right|\le91R .
\]
Appendix~\ref{app:emp} proves this estimate with the stated constant. Adding
the terms and dividing by $n$ gives $(320+81+182)R/n=583\sqrt{d/n}$.
Since the empirical losses take values in the finite set $\{0,1/n,\ldots,1\}$, the infimum is attained. Each output name realizes its trace, so the calculation is its actual
empirical error. This proves \eqref{eq:accuracy}.

\section{Deletion stability}\label{sec:stability}
For laws on a finite common alphabet use
$H^2(p,r)=\sum_h(\sqrt{p_h}-\sqrt{r_h})^2$.
Thus $H$ is an ordinary Euclidean distance between square-root vectors.
For a labeled block $T$ of size $N=n+1$, let $q(U)$ be its optimizing
law on the canonical names. Let $q_{-i}(U_{-i})$ be the optimizing law on
deleting position $i$, also expressed on the full names. The learner
averages over the weights; these conditional laws allow us to couple the
two experiments by using the same surviving weights.

\begin{proposition}[Total deletion cost]\label{prop:deletion}
For every such $T$,
\[
 \sum_{i=1}^N H^2\bigl(\E_U q(U),\E_U q_{-i}(U_{-i})\bigr)\le4928d.
\]
\end{proposition}
\begin{proof}
Deleting a position changes its loss, the collection of feature subsets
used by the entropy penalty, and possibly the name of a trace. We account for
these changes in that order. Write $r_i$ for the optimizer after setting
$U_i=0$ while retaining all features. There are two estimates:
\begin{align}
 \sum_i\E_U H^2(q,r_i)&\le64d, \label{eq:lossremove}\\
 \sum_i H^2(q,q_{-i})&\le576d+68\sum_i H^2(q,r_i).
 \label{eq:fiberremove}
\end{align}
The second is pointwise in $U$. We explain the mechanism here; Appendices
\ref{app:projection-bounds} and \ref{app:deletion} prove the estimates.

At a positive law $q$, conditional expectation given the predictions on $M$
is an orthogonal projection $P_M$ in $L^2(q)$. Its average over $m$-subsets is
$K_m$. On mean-zero functions the Hessian of the entropy penalty in score
coordinates is $A_q=\tau I+\sum_j a_j K_{m_j}$. If only $U_i$ changes, the
score $s_i=(\partial_i q)/q$ solves
$A_qs_i=-(\ell_i-\E_q\ell_i)$. Thus the total squared speed of the
square-root law is controlled by the sum of squared inverse-Hessian
responses to the individual centered losses.

A prediction already determined by $m$ observed
coordinates has zero residual variance. For the remaining predictions the
one-inclusion graph \cite{hlw} with weights given by the trace probabilities
bounds the average residual variance by
$d/(m+1)$. Averaged projections at very different subset sizes also overlap
weakly in their unrecovered directions. Combining these facts at the
geometrically spaced $m_j$ gives $\sum_i\|s_i\|_q^2\le64d$ for every
positive optimizing law. Appendix~\ref{app:deletion} integrates the speed
along each coordinate, then averages over the independent weights to obtain
\eqref{eq:lossremove}.

For \eqref{eq:fiberremove}, remove $i$ from the entropy restrictions and
project full traces onto the surviving positions. At most two full traces
can merge, since just one binary prediction was forgotten. Only such a pair
can force a change of output name. The total mass of the lighter endpoints
over all pairs and deletions is at most $d$ by the same one-inclusion bound.
Appendix~\ref{app:deletion} bounds the changes arising from the new restrictions, the merged traces, and
the adjustment to the true deleted optimizer, yielding \eqref{eq:fiberremove}.

Finally $H^2(\E p,\E r)\le\E H^2(p,r)$ by Cauchy--Schwarz applied to
$\E\sqrt{p_h r_h}$ on each atom. Take expectations in
\eqref{eq:fiberremove} and use \eqref{eq:lossremove}. The resulting
constant is $576+68\cdot64=4928$.
\end{proof}

\section{Conditional mutual information}\label{sec:info}
We pass from Proposition~\ref{prop:deletion} to the fixed-pool information
without a measurable-space structure. Fix an ordered pool $t$ and put
$Y=A(t_S)$. The chain rule and independence of the bits give
\begin{equation}
 I(S;Y)\le \sum_{i=1}^n I(S_i;Y\mid S_{-i}).
 \label{eq:tensor}
\end{equation}
In the chain-rule term for bit $i$, $H(S_i\mid S_{<i})=\log2$.
Conditioning on the remaining bits as well as $Y$ can only decrease
$H(S_i\mid Y,S_{<i})$, proving the inequality in this direction.

Fix the other bits. There are now just two possible input samples. Let their
output laws be $p,r$. The information of this one fair bit is
$\JS(p,r)=[\KL(p\|v)+\KL(r\|v)]/2$, with $v=(p+r)/2$. From
$\log x\le x-1$ on the positive terms and
$(\sqrt p+\sqrt r)^2\le2(p+r)$ coordinatewise,
\[
 \JS(p,r)\le\tfrac12\sum_{h:p_h+r_h>0}\frac{(p_h-r_h)^2}{p_h+r_h}
             \le H^2(p,r).
\]
Form an ordered block $T$ by listing the chosen observation in each other
row and listing \emph{both} observations in row $i$, consecutively in
positions $i,i+1$. Removing either of these two positions gives one of
the training samples, in its original row order. Let $Q$ be the full-block
learner law computed with the schedule for $n$, and let $c_j(T)=H^2(Q,Q_{-j})$, using weight-averaged laws.
The triangle inequality for square-root vectors gives
$H^2(p,r)\le2(c_i(T)+c_{i+1}(T))$.

Now hold the $2n$ observations fixed and shuffle their \emph{positions}.
For any prescribed row and values of the other bits, the ordered $N$-block
just constructed has the same law as the first $N$ positions of a uniform
permutation. Indeed every ordered injection of $N$ positions has exactly
$(2n-N)!$ extensions to a permutation. This remains true when observation
values coincide. Write $u_j=\E_T c_j(T)$ for that common block law.
Summing \eqref{eq:tensor} after the shuffle gives
\[
 \E_\pi F_A(t\circ\pi)\le 2\sum_{i=1}^n(u_i+u_{i+1})
           \le 4\sum_{j=1}^{n+1} u_j
           \le4\cdot4928d=19{,}712d.
\]
The middle inequality counts how often one block position occurs in adjacent
pairs. Coherence embeds both deletion alphabets into the full-block alphabet
of named hypotheses without identifying distinct names. This proves
Theorem~\ref{thm:pool}.

\section{Measurability}\label{sec:scope}
We first establish integrability of the scalar statistic under a given law.
We then give sufficient conditions on the class for a distribution-independent
measurable kernel.

\paragraph{The iid transfer.}
For fixed $n,d$, $F_A$ takes finitely many possible values, even as the
features and their canonical names vary over infinite spaces. To see this,
record the labels and full trace family of the $2n$ positions. For each
selector $s$, record the training traces and the full-pool trace of the
representative chosen for each training trace. There are finitely many
such records. The optimization depends only on the finite training family,
its labels and the weights. The record tells us which returned names from
different selections are the same: each of them is the canonical name of
its recorded full-pool trace. Consequently the record determines every
summand of \eqref{eq:fixed}.
If $F_A$ is almost everywhere measurable under the iid product law, its
finite range gives integrability. For each position permutation $\pi$,
$Z\circ\pi$ has the law of $Z$. Integrating \eqref{eq:pool} and interchanging
a finite average and the integral gives $\E_Z F_A(Z)\le19{,}712d$.

\begin{definition}[Measurable canonical representatives]\label{def:selection}
Give $\XX$ and $\HH$ measurable structures, and fix the order used by $g$.
Require the following events to be measurable in every finite tuple $x$:
\begin{enumerate}
\item A specified binary trace $v$ is realized on $x$.
\item For a fixed map of positions $e:K\to I$, a specified trace $v$
is realized on $x\circ e$ and its representative $g_{x\circ e}(v)$ has
a specified trace on all of $x$.
\item A specified trace $v$ is realized on $x$ and $g_x(v)$ belongs to a
specified measurable output set $B\subseteq\HH$.
\end{enumerate}
The assertions involving a representative are made only on the event
that the trace is realized.
\end{definition}

The first two conditions make every cell of the preceding finite record
measurable. For the kernel itself, fix an output event $B$. On a given
training trace family and label pattern the weight-averaged probability of
each trace is a constant determined by \eqref{eq:objective}. Sum those
constants times the indicators from condition 3 over the finitely many
traces. This gives the probability of $B$ and makes it measurable.
Under this kernel, conditional on $Z=t$, choose $s$ uniformly and then
draw from $A(t_s)$. This is a measurable joint kernel for $(S,Y)$ given
the pool: on a selector value $s$ and an output event $B$ it assigns
mass $2^{-n} A(t_s)(B)$. Thus it supplies the conditional law appearing
in the authors' definition. With measurable output singletons the finite
set of names on each pool is
measurably distinguishable, so its discrete KL agrees with KL of the joint
measure relative to the product of its marginals. Coarsening the output
sigma-algebra decreases KL by data processing. These facts prove
Theorem~\ref{thm:iid}.

\subsection{A sufficient condition from ordered empirical risk minimization}

\begin{proposition}\label{prop:ordered-erm}
Suppose evaluation $(h,x)\mapsto h(x)$ is jointly measurable. Fix a well-order
of $\HH$. If the first empirical risk minimizer in this order is a measurable
$\HH$-valued function of the labeled sample at every size, then
Definition~\ref{def:selection} holds for that order.
\end{proposition}
\begin{proof}
Given a tuple $x$ and a prescribed binary vector $v$, run the minimizer on the
sample obtained by labeling $x$ with $v$. Denote its output by $a_v(x)$.
This is a measurable function of $x$, since the labels are fixed.
The vector $v$ is realizable precisely when $a_v(x)$ fits all these labels:
if any classifier has zero mistakes, every minimizer has zero mistakes.
Joint evaluation makes this finite test measurable. On its success event the
tie-breaking rule gives $a_v(x)=g_x(v)$, since the hypotheses of zero loss
are exactly the ones realizing $v$.

For a subtuple, run the same procedure there and evaluate its output on all
positions of the larger tuple. Requiring the subtuple test to succeed and
these further predictions to equal a given vector gives the second event
in Definition~\ref{def:selection}. Requiring the full-tuple test to succeed
and $a_v(x)$ to lie in a prescribed measurable output set gives the third.
\end{proof}

The minimizer in this proposition is the one constructed in the proof of
Lemma~4.15 of \cite{framework}. Its role here is to supply representatives.
The two rates concern the randomized learner of Section~\ref{sec:learner}.

\subsection{Countable domains and countable trace cores}
If $\HH$ is countable and each predictor is measurable, the event that
a particular $h$ is selected is the intersection of its trace event with
the complements of the trace events of all earlier hypotheses. This
is a countable intersection of measurable sets. Taking countable unions
over candidate $h$ proves all three clauses, for any prescribed order.
If instead $\XX$ is countable and
singletons are measurable, then every finite power of $\XX$ is countable
with measurable singletons. Every subset of that power is measurable.
This proves all three clauses for arbitrary $\HH$, including thresholds
with real cut points acting on rational inputs. Countable support of a
particular $D$ similarly gives measurability of the scalar statistic on its
full-measure countable pool support when feature singletons are measurable.
Here the full-measure set can depend on $D$.

For an uncountable class it suffices to have a countable
subclass $C$ of measurable predictors realizing every finite trace of
$\HH$. Place a fixed enumeration of $C$ before the rest of a well-order.
Every selected name lies in $C$. The preceding countable unions and
intersections can therefore be taken inside $C$. This preserves coherence
across all deletions, and $C$ and $\HH$ have the same empirical optimum
on each sample. For example, rational cut points realize every finite
trace of the real thresholds $h_a(x)=\ind\{x\ge a\}$ on $\mathbb R$.
When both predictions occur, choose a rational strictly above the largest
excluded observation and at or below the smallest included observation;
cuts below or above all observations give the two constant traces.
More generally a separable parameter space with continuous binary
evaluation in the parameter at each instance has this property if its
predictors are measurable. Each realized finite pattern defines a nonempty
open set of parameters, so a fixed countable dense set meets it. A second
sufficient condition is that every classifier is the pointwise limit of a
sequence from one countable measurable subclass. Binary convergence gives
eventual equality at each point, hence simultaneous equality on a finite tuple.

\subsection{An obstruction to Borel selection}
Let $B\subset[0,1]$ be non-Borel. Take the zero classifier and the singleton
indicators $\ind_{\{b\}}$ for $b\in B$. Singletons are Borel, the class
shatters one point of $B$, and no classifier labels two distinct points
positive. Yet the one-point event that some hypothesis predicts 1 at $x$
is exactly $B$. Condition 1 fails for every order.

The same example constrains other proper learners. Suppose the output rule
is a Borel kernel and evaluation is jointly measurable. Give the learner $n$ copies of $(x,1)$ and write $p(x)$ for
the probability of returning a hypothesis that predicts 1 at $x$.
This is a Borel function: integrating a jointly measurable bounded
function against a kernel preserves measurability, first for measurable
rectangles and then by the monotone-class theorem. In the singleton
example joint measurability holds on the disjoint union of a point and
$B$ with its trace sigma-algebra, since evaluation is a diagonal equality.
For $x\notin B$ properness forces $p(x)=0$. For $x\in B$, the empirical
optimum is zero, so \eqref{eq:conj} forces $p(x)\ge1-c/\sqrt n$.
Fix $c$ and take an integer $n>\max\{1,c^2\}$. This makes $B=\{x:p(x)>0\}$, a contradiction.

\begin{proposition}
There are standard Borel instance and hypothesis spaces with jointly Borel
evaluation and VC dimension one such that, for every $c\ge0$ and integer
$n>\max\{1,c^2\}$, no proper Borel Markov learner satisfies the empirical
inequality in \eqref{eq:conj} on every sample.
\end{proposition}
\begin{proof}
Use the classical fact
that there exists an analytic non-Borel subset $B$ of Baire space
$P=\mathbb N^{\mathbb N}$ and that $B$ is the image of $P$ under a
continuous map $f$ \cite[Def.~14.1 and Thm.~14.2]{kechris}.
Take $\XX$ to be two disjoint copies of $P$. In addition to zero, include
one hypothesis $h_u$ for each $u\in P$, with positive set
$\{(0,f(u)),(1,u)\}$. Give the parameter space $\{*\}\sqcup P$ its standard Borel structure
and identify $*$ with the zero hypothesis. The parameter $u$ is identifiable
by its positive point in copy 1, so this also gives a measurable structure
on the class of distinct functions. Evaluation is jointly Borel, as it uses the graph of $f$
and the diagonal. No two points from the same copy can both be positive.
For a pair meeting both copies, positivity at the point in copy 1 determines
the entire hypothesis, preventing the two patterns that agree there with
value 1 and differ in copy 0. Thus the VC dimension is one. Repeating
the instance $(0,x)$ with label 1 detects precisely $f(P)=B$ by the same
argument, contradicting Borel measurability.
\end{proof}
This proposition concerns Borel measurability;
analytic sets are universally measurable \cite[Thm.~21.10]{kechris},
so this example does not give the same obstruction for universally measurable kernels.

\paragraph{Relation to Conjecture 6.}
The problem paper uses random variables in Definition~3. Its companion
paper defines CMI as the expectation of conditional mutual informations and
uses a well-order for global tie-breaking \cite[Defs.~2.6, 2.8 and
Lemmas~4.15, 4.17]{framework}. Proposition~\ref{prop:ordered-erm} shows that, under joint
measurability of evaluation, measurability of this particular minimizer
suffices for both rates. The well-order construction proves consistency of
its choices but gives no measurability assertion. Indeed, the preceding
VC-one example also rules out an everywhere Borel exact minimizer: such a
minimizer would satisfy the empirical inequality with zero excess.
Conjecture~6 leaves the measurable-space convention for arbitrary classes
unspecified. The population theorem therefore remains conditional in that
generality.

\appendix
\section{The empirical-process estimate with constants}\label{app:emp}
For completeness we derive the bound used in \eqref{eq:accbudget}.
For a nonempty binary family on $m$ positions, split at one position into
two projected layers. Their union has VC dimension at most $d$ and their
intersection has dimension at most $d-1$: shattering a set in the latter
would shatter it together with the split coordinate in the original family.
The sum of the layer sizes equals the sum of their union and intersection
sizes. Induction using Pascal's identity gives
$|V|\le\sum_{k=0}^d\binom{m}{k}$, with at most one trace at dimension zero
and the trivial count when $d\ge m$. For $1\le d\le m$, multiply the sum
by $(d/m)^d$, bound each term by $\binom{m}{k}(d/m)^k$, and apply
$(1+d/m)^m\le e^d$. This proves the trace count in Section~\ref{sec:accuracy}.

Measure Hamming distance as the \emph{fraction} of disagreeing positions.
Suppose $p$ traces are separated by at least $\delta\in(0,1]$.
Draw $k=d+\lceil2\log(p)/\delta\rceil+1$ positions with replacement.
The probability that any pair agrees on them all is at most
$p^2e^{-k\delta}<1$. Hence the packed traces restrict injectively to some drawn tuple.
Repeated positions cannot increase its VC dimension, by the argument in
Section~\ref{sec:accuracy}. Applying the trace count to this tuple with $u=(\log p)/d$, we have
\[
 u\le1+\log(k/d),\qquad k/d\le3(1+u)/\delta.
\]
Indeed $k/d\le1+2u/\delta+2/d$. Since
$\log(1+u)\le\log2+(u-1)/2$, rearranging gives
$\log p\le d(1+2\log(6/\delta))$.

Let $r_s=2^{-s}$, $s\ge0$, and take a largest $r_s^2$-separated family
$G_s$ of traces. Every trace lies within Hamming distance $r_s^2$ of
$G_s$, since otherwise it could be added. At a sufficiently fine scale
$G_s$ is the full family, since different traces differ on at least one
of $n$ positions. The packing estimate gives
$\log(2|G_s|)\le12d(s+1)$. Join each member of $G_{s+1}$ to a chosen
parent in $G_s$ within distance $r_s^2$.

Fix amplitudes $|\alpha_i|\le1$ and let $X(v)=\sum_i\varepsilon_i\alpha_i\ell_i(v)$
for independent fair signs $\varepsilon_i$. A signed sum with squared
coefficient norm at most $nr^2$ has exponential moment at most
$e^{\lambda^2nr^2/2}$, by $\cosh x\le e^{x^2/2}$. Thus for $m$ such
sums, Jensen's inequality and summing the positive and negative exponentials
give expected maximum absolute value at most
$\log(2m)/\lambda+\lambda nr^2/2$. Set $\lambda=d/(Rr)$, $R=\sqrt{nd}$.
If $\log(2m)\le dL$, the bound becomes $(L+1/2)Rr$.
At scale zero use $r=1,L=12$, for a bound $13R$.
For the increments $X(v)-X(\mathrm{parent}(v))$ at level $s+1$,
only disagreements contribute, so their squared coefficient norms are
at most $nr_s^2$. There are at most $|G_{s+1}|$ such increments.
Their expected maximum is at most $13(s+2)R2^{-s}$.
Telescope along the parents, then sum the maxima. The geometric series
$\sum_{s\ge0}(s+2)2^{-s}=6$ gives $(13+13\cdot6)R=91R$.
Conditioning on the absolute values of $U_i-1$ gives the desired bound.

\section{Projection estimates and the score bound}
\label{app:projection-bounds}

Let $[N]=\{1,\ldots,N\}$, let $F\subseteq\{0,1\}^{N}$ be a nonempty trace family,
and let $q$ be a positive probability law on $F$. Write
\[
  \langle f,g\rangle_q=\sum_{\xi\in F}q(\xi)f(\xi)g(\xi),
  \qquad \|f\|_q^2=\langle f,f\rangle_q.
\]
For $M\subseteq[N]$, let $U_M$ be the subspace of functions that are constant on
fibers of $\xi\mapsto\xi|_M$, and let $P_M$ be its orthogonal projection.
On each fiber, the value of $P_M f$ is the $q$-weighted mean of $f$. Indeed, the
residual from this mean has weighted sum zero on every fiber, so it is orthogonal
to $U_M$. The map $f\mapsto(\sqrt{q(\xi)}f(\xi))_{\xi\in F}$ identifies this
weighted space with ordinary Euclidean space.

For $0\leq m\leq N$, define $K_m=\mathbb E_{|M|=m}P_M$, with $M$ uniform.
If $M\subseteq V$, then $U_M\subseteq U_V$. The difference $P_V-P_M$
projects onto $U_V\cap U_M^\perp$, giving $P_M\preceq P_V$ in the
positive-semidefinite order. Apply a uniform permutation of $[N]$ to any fixed
nested pair of sets of sizes $m$ and $m'$. Both marginals are uniform, giving
$0\preceq K_m\preceq K_{m'}\preceq I$ whenever $m\leq m'$. The following
argument applies to any nested family of subspaces of a finite-dimensional
real Hilbert space.

\begin{lemma}[Overlap of averaged projections]
\label{lem:projection-overlap}
If $k,m,M$ are integers with $k\geq1$, $m\geq0$, and $km\leq M\leq N$,
then for every $y$,
\begin{equation}
  \label{eq:projection-overlap-form}
  \big\langle(I-K_M)K_my,K_my\big\rangle
  \leq \frac{1}{k}\langle K_my,y\rangle.
\end{equation}
In particular, $\|(I-K_M)^{1/2}K_m^{1/2}\|_{\mathrm{op}}\leq k^{-1/2}$.
\end{lemma}

\begin{proof}
Draw $J_1,\ldots,J_k$ independently and uniformly among the $m$-subsets
of $[N]$. For each tuple, choose any $M$-subset $V_0$ containing its union,
which has size at most $km$. Independently draw a uniform permutation $\pi$
of $[N]$ and set $V=\pi(V_0)$ and $J'_s=\pi(J_s)$. Conditional on the
original tuple, $V$ is uniform among the $M$-subsets. Conditional on $\pi$,
the $J'_s$ are independent and uniform, and this joint law is the same for
every $\pi$. Also $J'_s\subseteq V$ for every $s$.

Put $x=K_my$ and $Y_s=P_{J'_s}y$. The average $k^{-1}\sum_s Y_s$ lies in
$U_V$. The best-approximation property of $P_V$ gives
\[
  \langle(I-K_M)x,x\rangle
  =\mathbb E\|x-P_Vx\|^2
  \leq\mathbb E\left\|x-\frac1k\sum_{s=1}^kY_s\right\|^2.
\]
The $Y_s$ are independent with mean $x$, so the cross terms between
$Y_s-x$ and $Y_t-x$ for $s\ne t$ have mean zero. Since
$\|P_Jy\|^2=\langle P_Jy,y\rangle$, the right side equals
$k^{-1}(\langle K_my,y\rangle-\|K_my\|^2)$. This proves
\eqref{eq:projection-overlap-form}.

Write $R=K_m^{1/2}$ and take any $v$. Let $w$ be the orthogonal projection
of $v$ onto $\operatorname{ran}R$, and choose $y$ with $Ry=w$. The symmetric
operator $R$ has $\ker R=(\operatorname{ran}R)^\perp$, so $K_my=Rv$ and
$\langle K_my,y\rangle=\|w\|^2\leq\|v\|^2$. Applying
\eqref{eq:projection-overlap-form} to $y$ proves the operator-norm assertion.
\end{proof}

For the rest of the appendix, take $1\leq d\leq n$, $N=n+1$,
$t=n/d$, and $\rho=21/20$. Let $J$ be the largest nonnegative integer with
$d16^J\leq n$. Set
\[
  m_j=d4^j,\qquad b_j=\sqrt t\,\rho^{-j},\qquad
  a_j=b_j-b_{j+1},\qquad \tau=b_{J+1}.
\]
In particular, $m_J\leq n<N$. Define the positive operators
\[
  B_0=K_{m_0},\qquad B_j=K_{m_j}-K_{m_{j-1}}\ (1\leq j\leq J),
  \qquad B_{J+1}=I-K_{m_J}.
\]
They sum to $I$. Collecting the coefficients of each $K_{m_j}$ gives
\[
  A=\sum_{j=0}^{J+1} b_jB_j=\tau I+\sum_{j=0}^{J}a_jK_{m_j}.
\]

\begin{lemma}[Multiscale inverse estimate]
\label{lem:multiscale-inverse}
The operator $A$ is invertible, and for every $x$,
\begin{equation}
  \label{eq:multiscale-inverse}
  \|A^{-1}x\|^2
  \leq 4\sum_{j=0}^{J+1}b_j^{-2}\langle B_jx,x\rangle.
\end{equation}
\end{lemma}

\begin{proof}
First consider positive operators $P\preceq K$ and $Q\preceq L$, with
$K$ positive, and suppose that
$\langle L Kz,Kz\rangle\leq c^2\langle Kz,z\rangle$ for every $z$ and some
$c\geq0$. For any $y$ put $\gamma=\langle KQy,Qy\rangle\geq0$.
Cauchy--Schwarz for a positive form, obtained by applying ordinary
Cauchy--Schwarz to the square root of its operator, gives
\[
  \gamma^2\leq\langle Qy,y\rangle\langle QKQy,KQy\rangle
  \leq c^2\langle Qy,y\rangle\gamma.
\]
Thus $\gamma\leq c^2\langle Qy,y\rangle$, including the case $\gamma=0$.
Applying Cauchy--Schwarz to the form for $P$ and using $P\preceq K$ yields
\begin{equation}
  \label{eq:projection-band-pair}
  |\langle Px,Qy\rangle|^2
  \leq c^2\langle Px,x\rangle\langle Qy,y\rangle.
\end{equation}
For $i<j$, take $P=B_i$, $Q=B_j$, $K=K_{m_i}$ and
$L=I-K_{m_{j-1}}$. Here $P\preceq K$, $Q\preceq L$, and
$m_{j-1}=4^{j-i-1}m_i$, so Lemma~\ref{lem:projection-overlap} gives
$c=2^{-(j-i-1)}$. Swapping the two bands handles $i>j$.
On the diagonal, use $K=L=I$ and $c=1$. Consequently
\eqref{eq:projection-band-pair} holds for $P=B_i$, $Q=B_j$ with
$c=c_{ij}=2^{-\max\{|i-j|-1,0\}}$.

Set $C_0=\sum_i b_i^{-1}B_i$. The bands are self-adjoint and
$\sum_{i,j}B_iB_j=(\sum_iB_i)(\sum_jB_j)=I$, so
\[
  \langle(C_0A-I)x,y\rangle
  =\sum_{i,j}(b_j/b_i-1)\langle B_jx,B_iy\rangle.
\]
For each fixed row or column, at most two indices occur at separation
$h\geq1$. At that separation,
$|b_j/b_i-1|\leq\rho^h-1$. Thus each row and column of
$w_{ij}=|b_j/b_i-1|c_{ij}$ sums to at most
\[
  r=2\sum_{h\geq1}(\rho^h-1)2^{-(h-1)}
   =2\left(\frac{\rho}{1-\rho/2}-2\right)
   =\frac{8}{19}<\frac12.
\]
For $x,y\ne0$, put
$u_j=\langle B_jx,x\rangle^{1/2}/\|x\|$ and
$v_i=\langle B_iy,y\rangle^{1/2}/\|y\|$. Both squared sequences sum to
one. The preceding bilinear identity and \eqref{eq:projection-band-pair}
bound its absolute value by
\[
  \|x\|\|y\|\sum_{i,j}w_{ij}u_jv_i
  \leq \frac{\|x\|\|y\|}{2}\sum_{i,j}w_{ij}(u_j^2+v_i^2)
  \leq r\|x\|\|y\|.
\]
The inequality also holds if either vector vanishes. It follows that
$\|C_0A-I\|_{\mathrm{op}}\leq 8/19$.

Since $A\succeq\tau I$ and $\tau>0$, the finite-dimensional operator
$A$ is invertible. Taking $z=A^{-1}x$, the error bound gives
$\|z\|\leq\|C_0x\|+\tfrac12\|z\|$. Finally, positive-form Cauchy--Schwarz
and the scalar Cauchy--Schwarz inequality give
\[
  |\langle C_0x,y\rangle|
  \leq
  \Big(\sum_j b_j^{-2}\langle B_jx,x\rangle\Big)^{1/2}
  \Big(\sum_j\langle B_jy,y\rangle\Big)^{1/2}.
\]
The second factor is $\|y\|$. Taking the supremum over unit $y$ and
squaring $\|z\|\leq2\|C_0x\|$ proves \eqref{eq:multiscale-inverse}.
\end{proof}

\begin{lemma}[Weighted one-inclusion bound]
\label{lem:weighted-one-inclusion}
Let $G$ be a finite binary trace family of VC dimension at most $d$.
Its one-inclusion graph has one edge per unordered pair of traces that
differ on a single coordinate. For any nonnegative vertex weights $w$,
\[
  \sum_{\{u,v\}\in E(G)}\min\{w(u),w(v)\}
  \leq d\sum_{u\in G}w(u).
\]
\end{lemma}

\begin{proof}
First prove $|E(G)|\leq d|G|$ by induction on the number of coordinates,
simultaneously for every $d$. With no coordinates there are no edges.
For $d=0$, any edge would shatter its differing coordinate.
For $d\geq1$, split at the last coordinate into projected layers $G_0,G_1$.
The union $U=G_0\cup G_1$ has dimension at most $d$. If the intersection
$V=G_0\cap G_1$ shatters a set, $G$ shatters that set together with the
last coordinate, so $V$ has dimension at most $d-1$.
For each potential edge in the remaining coordinates, the number of layers
in which both endpoints occur is at most the number of families among
$U,V$ in which both occur. There are $|V|$ edges across the layers.
The induction bound gives
\[
  |E(G)|\leq |E(U)|+|E(V)|+|V|
     \leq d(|U|+|V|)=d|G|.
\]

Let $0=\gamma_0<\gamma_1<\cdots<\gamma_L$ list zero and the distinct
positive weights and set $G_\ell=\{u:w(u)\geq\gamma_\ell\}$. The weight
of any edge is the sum of $\gamma_\ell-\gamma_{\ell-1}$ over levels with
both endpoints in $G_\ell$. Each $G_\ell$ has VC dimension at most $d$.
Summing the unweighted bound with these coefficients gives
\[
  \sum_{\ell=1}^L(\gamma_\ell-\gamma_{\ell-1})|E(G_\ell)|
  \leq d\sum_{\ell=1}^L(\gamma_\ell-\gamma_{\ell-1})|G_\ell|
  =d\sum_{u\in G}w(u).
\]
\end{proof}

For every coordinate $i$ write $f_i(\xi)=\xi_i-\mathbb E_q\xi_i$.

\begin{lemma}[Average conditional variance]
\label{lem:conditional-variance}
If $F$ has VC dimension at most $d$ and $0\leq m<N$, then
\begin{equation}
  \label{eq:averaged-conditional-variance}
  \frac1N\sum_{i=1}^N\langle(I-K_m)f_i,f_i\rangle_q
  \leq\frac{d}{m+1}.
\end{equation}
\end{lemma}

\begin{proof}
For fixed $i,M$, the quadratic form for $I-P_M$ is the mean
conditional variance of $\xi_i$ given $\xi|_M$. It vanishes if $i\in M$.
Draw $M$ uniformly among the $m$-subsets and $i$ uniformly from $[N]$,
independently. Conditional on $i\notin M$ and on $T=M\cup\{i\}$, the
position $i$ is uniform in $T$.

Push $q$ forward to a law $p$ on the projected family $F_T$. A fiber
of the projection from $T$ to $T\setminus\{i\}$ has at most two traces.
A two-element fiber of masses $p_u,p_v$ contributes
$p_up_v/(p_u+p_v)\leq\min\{p_u,p_v\}$ to the expected variance.
A singleton fiber contributes zero. Summing over the omitted positions
in $T$ visits each one-inclusion edge exactly once. Restriction cannot
increase VC dimension, so Lemma~\ref{lem:weighted-one-inclusion} bounds
this sum by $d\sum_{v\in F_T}p_v=d$. Averaging over the $m+1$ omissions
and including the zero contributions when $i\in M$ proves
\eqref{eq:averaged-conditional-variance}.
\end{proof}

\begin{proposition}[Score energy for the geometric schedule]
\label{prop:score-energy}
For every positive probability law $q$ on $F$ with $\operatorname{VCdim}(F)\leq d$,
\begin{equation}
  \label{eq:score-energy}
  \sum_{i=1}^N\|A^{-1} f_i\|_q^2\leq 64d.
\end{equation}
For the entropy-regularized optimizer at fixed weights, the same bound
holds for $\sum_i\|s_i\|_q^2$, where $s_i(\xi)=\dot q_i(\xi)/q(\xi)$
and $\dot q_i$ is the derivative in the $i$th loss weight.
\end{proposition}

\begin{proof}
Since $B_0\preceq I$ and a binary variable has variance at most $1/4$,
we have
$\sum_i\langle B_0f_i,f_i\rangle_q\leq N/4$.
For $1\leq j\leq J+1$ we have $B_j\preceq I-K_{m_{j-1}}$.
Lemma~\ref{lem:conditional-variance} therefore gives
\[
  \sum_i\langle B_jf_i,f_i\rangle_q
  \leq\frac{Nd}{m_{j-1}+1}\leq N4^{1-j}.
\]
The last upper bound also holds at $j=0$. Applying
\eqref{eq:multiscale-inverse} and using $b_j^{-2}=t^{-1}\rho^{2j}$ gives
\[
  \sum_i\|A^{-1}f_i\|_q^2
  \leq\frac{16N}{t}\sum_{j=0}^{J+1}\left(\frac{\rho^2}{4}\right)^j
  \leq\frac{32N}{t}\leq64d.
\]
Here $\rho^2/4=441/1600<1/2$ and $N\leq 2n$.

For an optimizing law, Lemma~\ref{lem:deletion-regularity} applies to
variation of just the $i$th weight. The score has mean zero since
$q$ has total mass one; each $P_M$ preserves means, so $As_i$ also
has mean zero. Testing the tangent equation against every mean-zero
function therefore gives
$A s_i=-(\ell_i-\mathbb E_q\ell_i)=\pm f_i$, since the label is binary.
Substitution into \eqref{eq:score-energy} proves the assertion for the
scores; the derivative and its endpoint interpretation are justified in
Appendix~\ref{app:deletion}.
\end{proof}

\section{Optimizer regularity and deletion estimates}
\label{app:deletion}

Fix integers $1\le d\le n$ and a labeled block
$T=((x_i,y_i))_{i\in I}$ with $|I|=N=n+1$. Let $G$ contain the least
hypothesis for each distinct binary prediction vector on the block, using the
same well-order as the learner. All laws compared below live on $G$, with zeros
where necessary. Write $R_M$ for the projection of a trace onto the positions
$M\subseteq I$. For a mass vector $p$ on $G$, write $p_M=(R_M)_\#p$ for the
projected vector. Use only realized cells as the codomain of $R_M$. Let
$\Phi(p)=\sum_h p(h)\log p(h)$, with $0\log 0=0$, and set
\[
  H^2(p,r)=\sum_{h\in G}\bigl(\sqrt{p(h)}-\sqrt{r(h)}\bigr)^2.
\]
In particular, this convention for squared Hellinger distance has no factor of
$1/2$.

The schedule is held fixed at the training size $n$ throughout the deletion
argument. Put $t=n/d$, $\rho=21/20$, and let $J$ be the largest integer with
$d16^J\le n$. For $0\le j\le J$ let $m_j=d4^j$ and $a_j=b_j-b_{j+1}$, where
$b_j=\sqrt t\,\rho^{-j}$ for $0\le j\le J+1$. Set $\tau=b_{J+1}>0$. The weight
of a restriction at level $j$ is $a_j/\binom N{m_j}$ on the full block. With
$\ell_i(h)=\mathbf 1\{h(x_i)\ne y_i\}$ and $L_U=\sum_i U_i\ell_i$, the full
optimizer $q$ minimizes
\begin{equation}
  \mathcal F_U(p)=\sum_{h\in G}p(h)L_U(h)+\tau\Phi(p)
       +\sum_{j=0}^J a_j\,\mathbb E_{|M|=m_j}\Phi(p_M).
  \label{eq:deletion-objective}
\end{equation}
The subset expectations are uniform. The same formula on the surviving traces
and the surviving positions defines the actual optimizer after deletion. The
proof below uses two estimates from the projection and VC arguments. Write
$\mathcal E_i$ for the undirected edges of the full trace family whose endpoints
differ only at $i$. For any probability law $p$ on $G$, the weighted
one-inclusion estimate is
\begin{equation}
  \sum_i\sum_{\{h,h'\}\in\mathcal E_i}\min\{p(h),p(h')\}\le d.
  \label{eq:deletion-input-edges}
\end{equation}
For positive $p$, let $P_M$ be conditional expectation on the predictions on
$M$, acting on mean-zero functions in $L^2(p)$. With
$K_m=\mathbb E_{|M|=m}P_M$ and
$f_i(h)=h(x_i)-\mathbb E_p h(x_i)$, the spectral estimate used here is
\begin{equation}
  A(p)=\tau\mathrm{Id}+\sum_{j=0}^J a_jK_{m_j},
  \qquad
  \sum_i\|A(p)^{-1}f_i\|_{L^2(p)}^2\le 64d.
  \label{eq:deletion-input-spectrum}
\end{equation}
These estimates hold for the full trace family, including repeated features.
The empirical argument can use the same optimizer through the point-mass
comparison proved next.

\subsection{Finite optimization and the tangent equation}

Consider any objective of the form
$\mathcal F(p)=\langle p,L\rangle+\tau\Phi(p)+\sum_k c_k\Phi((R_k)_\#p)$,
where $G$ is finite and nonempty, each $R_k$ is a finite restriction map,
$\tau>0$, and $c_k\ge0$. Continuity of $x\log x$ on $[0,1]$ gives existence on
the compact simplex. Strict convexity of $x\log x$ on $[0,1]$ makes the
full-entropy term strictly convex on the simplex. The remaining entropy terms
are convex because pushforward is linear. Thus the minimizer is unique.

To see that it is positive, suppose a minimizing $p$ vanishes at $h_0$. Set
$p_\epsilon=(1-\epsilon)p+\epsilon\delta_{h_0}$ for $0<\epsilon<1$. Its full
negative entropy is
\[
  \Phi(p_\epsilon)=(1-\epsilon)\Phi(p)
       +(1-\epsilon)\log(1-\epsilon)+\epsilon\log\epsilon.
\]
The projected law of a point mass has zero entropy. By convexity of each
projected term and linearity of the loss,
\[
  \mathcal F(p_\epsilon)-\mathcal F(p)
  \le \epsilon\bigl(L(h_0)-\mathcal F(p)\bigr)
       +\tau\bigl((1-\epsilon)\log(1-\epsilon)
                   +\epsilon\log\epsilon\bigr).
\]
The right side is negative for sufficiently small $\epsilon>0$, since
$\tau>0$ and $\log\epsilon\to-\infty$. This proves positivity. For any
comparator $h_0$, minimization against $\delta_{h_0}$ also gives
\[
  \langle p,L\rangle\le L(h_0)-\tau\Phi(p)-\sum_k c_k\Phi((R_k)_\#p).
\]
If $|L(h)-L_0(h)|\le D$ uniformly, passing to $L_0$ on both sides costs at
most $2D$. Negative entropy on $k$ possible cells is at least $-\log k$ by
Jensen's inequality for $x\log x$ at the uniform mean of the $k$ masses.

For a positive law $p$ and a zero-mass tangent vector $v$, write $s=v/p$
coordinatewise. Then $\mathbb E_p s=0$. Write $\langle s,u\rangle_p=\sum_h p(h)s(h)u(h)$,
and let $P_k$ be conditional expectation given $R_k$. The gradient of a
projected negative entropy at coordinate $h$ is $1+\log p_k(R_k(h))$,
where $p_k=(R_k)_\#p$. The
Hessian of that term, on tangent directions $ps$ and $pu$, is
\begin{equation}
  \sum_c \frac{((R_k)_\#(ps))(c)((R_k)_\#(pu))(c)}{p_k(c)}
    =\langle P_k s,P_k u\rangle_{L^2(p)}.
  \label{eq:deletion-hessian}
\end{equation}
Indeed, $((R_k)_\#(ps))(c)/p_k(c)$ is the conditional mean of $s$ in the
cell $c$. The corresponding full-entropy Hessian is
$\langle s,u\rangle_{L^2(p)}$. In mass coordinates the full objective Hessian
$B$ therefore satisfies
\begin{equation}
  v^{\mathsf T}Bv\ \ge\ \tau\sum_h\frac{v(h)^2}{p(h)}.
  \label{eq:deletion-hessian-positive}
\end{equation}

\begin{lemma}[Regularity of the minimizer]
\label{lem:deletion-regularity}
Suppose $L$ and the coefficients $c_k$ vary affinely with a real parameter $w$,
with all $c_k(w)\ge0$ on the interval of interest and with $\tau>0$ fixed.
Then the minimizing law $p(w)$ is continuously differentiable on that
interval, using one-sided derivatives at its endpoints. If
$G_w(h)=\partial_w\!\left(\nabla\mathcal F_w\right)(p(w))(h)$ differentiates
the objective gradient while holding $p(w)$ fixed, the score
$s=\dot p(w)/p(w)$ satisfies, for every mean-zero $u$,
\begin{equation}
  \tau\langle s,u\rangle_p+\sum_k c_k(w)\langle P_k s,P_k u\rangle_p
    =-\langle G_w,u\rangle_p.
  \label{eq:deletion-tangent}
\end{equation}
\end{lemma}

\begin{proof}
At the positive optimum the gradient is constant in $h$, with common value
$\lambda$. Apply the implicit function theorem to
\[
  (w,\lambda,p)\longmapsto
      \left(\sum_h p(h)-1,\ \nabla\mathcal F_w(p)-\lambda\mathbf 1\right).
\]
All its terms are continuously differentiable near a positive optimum. For a
vector $(\dot\lambda,v)$ in the kernel of the derivative in $(\lambda,p)$,
we have $\sum_h v(h)=0$ and $Bv=\dot\lambda\mathbf 1$. Taking the Euclidean
inner product with $v$ gives $v^{\mathsf T}Bv=0$. By
\eqref{eq:deletion-hessian-positive}, $v=0$ and then $\dot\lambda=0$.
The derivative is a square finite-dimensional map and is therefore invertible.
The local solution remains positive. Wherever the coefficients are admissible,
it is the unique minimizer by convexity and the stationarity condition. This
also proves regularity up to the interval endpoints. Differentiating the
stationarity equation and pairing with $pu$ eliminates $\dot\lambda$.
Substituting \eqref{eq:deletion-hessian} gives \eqref{eq:deletion-tangent}.
\end{proof}

For any such path, coordinatewise differentiation of the square roots gives
\begin{equation}
  \left\|\frac{d}{dw}\sqrt{p(w)}\right\|_2^2
    =\frac14\sum_h\frac{\dot p(w,h)^2}{p(w,h)}
    =\frac14\|s\|_{L^2(p(w))}^2.
  \label{eq:deletion-speed}
\end{equation}
The fundamental theorem of calculus followed by Cauchy--Schwarz gives
\begin{equation}
  H^2(p(b),p(a))\le\frac{b-a}{4}\int_a^b
             \left\|\frac{\dot p(w)}{p(w)}\right\|_{L^2(p(w))}^2\,dw
  \qquad(a\le b).
  \label{eq:deletion-path}
\end{equation}

\subsection{Removing the loss of a row}

Let $r_i$ be the optimizer with $U_i$ set to zero and all restrictions and
traces unchanged. For fixed $U_{-i}$ let $p_i(w)$ be the optimizer with weight
$w$ at $i$. Lemma~\ref{lem:deletion-regularity} applies to this loss-only
path. Equation~\eqref{eq:deletion-tangent} on the mean-zero space is
\[
  A(p_i(w))\frac{\dot p_i(w)}{p_i(w)}
      =-\bigl(\ell_i-\mathbb E_{p_i(w)}\ell_i\bigr).
\]
The right side is $-f_i$ if $y_i=0$ and $f_i$ if $y_i=1$, with the centering
always taken at $p_i(w)$. Put $E_i(p)=\|A(p)^{-1}f_i\|_{L^2(p)}^2$. Taking
$a=0$ and $b=U_i\le2$ in \eqref{eq:deletion-path} yields
\begin{equation}
  H^2(q,r_i)\le\frac12\int_0^2 E_i(p_i(w))\,dw.
  \label{eq:deletion-loss-pointwise}
\end{equation}
The right side is the conditional mean for an independent
uniform replacement of $U_i$ on $[0,2]$. The same implicit function argument
with the full weight vector as parameter gives continuity of the optimizer
and of its partial derivatives. The square-root distances are bounded, and
the nonnegative score energies are bounded by \eqref{eq:deletion-input-spectrum}.
Consequently the averages can be taken under the product weight law. By
Fubini and \eqref{eq:deletion-input-spectrum},
\begin{equation}
  \mathbb E_U\sum_i H^2(q,r_i)\le\mathbb E_U\sum_i E_i(q)\le 64d.
  \label{eq:deletion-loss-budget}
\end{equation}

\subsection{Removing the row from the restriction scheme}

On the full support and with $U_i=0$, let $r_i'$ use only restrictions avoiding
$i$. Couple the old and new restrictions at size $m$ by drawing $M$ uniformly
among the $m$-subsets of $I$. Keep $M$ if it avoids $i$. Otherwise replace $i$
by a uniform element of $I\setminus M$. For a fixed target $M'$ avoiding $i$,
there is one unchanged source $M=M'$ and one source
$(M'\setminus\{j\})\cup\{i\}$ for each $j\in M'$. The target probability is
therefore
\[
  \frac{1}{\binom Nm}\left(1+\frac{m}{N-m}\right)
     =\frac{1}{\binom{N-1}{m}}.
\]
The probability of a change is $\binom{N-1}{m-1}/\binom Nm=m/N$. The
schedule ensures $1\le m\le n=N-1$, so both denominators are positive.

A changed pair consists of two refinements of a common restriction, each by
one binary prediction. At a positive law $p$, the gradient increment from
the common restriction to one refinement is the conditional log probability
of that binary prediction. If $c$ is the common cell and $b$ the binary
prediction in that cell, the two gradient constants $1$ cancel and give
$\log p(c,b)-\log p(c)=\log p(b\mid c)$. For $0<u\le1$, applying $\log x\le x-1$ at $x=u^{-1/2}$ gives
\[
  0\le -\sqrt u\log u
       \le 2(1-\sqrt u)\le2,
  \qquad u\log^2u\le4.
\]
At $u=0$ the squared-log term weighted by $u$ is zero by continuity. Summing
within a cell over the two binary values and averaging over the common cells
bounds the squared $L^2(p)$ norm of one increment by $8$. The difference of
two such increments consequently has squared norm at most $32$, by
$\|v-w\|^2\le2\|v\|^2+2\|w\|^2$.

Interpolate the two objectives on $[0,1]$ by assigning the old restriction
weights a factor $1-w$ and the new ones a factor $w$. The full-entropy
coefficient stays $\tau$. The forcing term $G_w$ in
\eqref{eq:deletion-tangent} is the weighted sum of the coupled gradient
differences. Unchanged pairs contribute zero. The total weight of changed
pairs across the levels is
\[
  W=\sum_{j=0}^J a_j\frac{m_j}{N}.
\]
If the changed gradients are $g_k$ with weights $w_k\ge0$, weighted
Cauchy--Schwarz gives
\[
 \Big\|\sum_k w_k g_k\Big\|_{L^2(p)}^2
 \le \Big(\sum_k w_k\Big)\sum_k w_k\|g_k\|_{L^2(p)}^2 .
\]
Thus $\|G_w\|_{L^2(p(w))}^2\le32W^2$ uniformly on the path. Testing
\eqref{eq:deletion-tangent} with $u=s$ gives
$\tau\|s\|^2\le-\langle G_w,s\rangle\le\|G_w\|\|s\|$. The case $s=0$ is
immediate. Otherwise divide by $\tau\|s\|$ and apply
\eqref{eq:deletion-path} on $[0,1]$ to obtain the bound
\begin{equation}
  H^2(r_i,r_i')\le B,
  \qquad B=8\left(\frac{W}{\tau}\right)^2.
  \label{eq:deletion-restriction-cost}
\end{equation}

To bound $NB$, extend the formula $m_k=d4^k$ through $k=J+1$ and write
$S_k=\sum_{j<k}a_jm_j$. The inequality $S_k\le b_km_k/4$ holds at $k=0$.
Using $b_k=\rho b_{k+1}$ and $m_{k+1}=4m_k$, the induction step gives
\[
  S_{k+1}\le\frac{b_km_k}{4}+(b_k-b_{k+1})m_k
     =\left(\frac54\rho-1\right)b_{k+1}m_k
     \le\frac{b_{k+1}m_{k+1}}4.
\]
Here $(5/4)\rho-1=5/16\le1$. Hence $S_{J+1}\le\tau m_J$.
The choice of $J$ gives $m_J^2=d^2 16^J\le nd$. With $W=S_{J+1}/N$,
\begin{equation}
  NB\le \frac{8m_J^2}{N}\le 8d.
  \label{eq:deletion-restriction-budget}
\end{equation}

\subsection{Collapsing redundant full traces}

Let $R_{-i}$ erase coordinate $i$. Each realized surviving trace has either
one or two extensions in $G$. Let $C_i$ be the set of surviving traces with
two extensions. The least hypothesis for a surviving trace is already a least
hypothesis for its full trace. An earlier hypothesis with the same full trace
would agree on every surviving coordinate and contradict minimality. Thus the
deleted canonical representative is an injective section $e_i$ of $R_{-i}$
with values in $G$.

Write $\bar p_i=(R_{-i})_\#r_i'$ and $\mu_i=\bar p_i(C_i)$. The zero-weight loss
and the restrictions used by $r_i'$ depend only on the surviving trace. At
fixed total mass $u$ on a two-element fiber, strict convexity of $x\log x$
minimizes the full negative entropy by assigning $u/2$ to each extension.
Thus $r_i'$ is the equal lift of $\bar p_i$. If $\mathcal F_i$ is the true
deleted objective on surviving traces, evaluation on the equal lift of any
law $\bar p$ gives
\begin{equation}
  \mathcal F_i(\bar p)-\tau(\log 2)\bar p(C_i).
  \label{eq:deletion-lift-objective}
\end{equation}
A two-element fiber contributes $u\log(u/2)=u\log u-u\log2$ to the full
negative entropy. A singleton contributes $u\log u$. The other terms agree
with the deleted objective because the surviving positions are merely
renumbered.

Let $\bar q_{-i}$ minimize $\mathcal F_i$, and let
$q_{-i}=(e_i)_\#\bar q_{-i}$ be the law on actual names in $G$. By
\eqref{eq:deletion-lift-objective}, $\bar p_i$ minimizes the perturbed coarse
objective. Comparing with $\bar q_{-i}$ gives
\begin{equation}
  \mathcal F_i(\bar p_i)-\mathcal F_i(\bar q_{-i})
      \le\tau(\log2)\bigl(\bar p_i(C_i)-\bar q_{-i}(C_i)\bigr)
      \le\tau(\log2)\mu_i.
  \label{eq:deletion-collapse-gap}
\end{equation}
The scalar remainder of $\phi(x)=x\log x$ at $y>0$ is
\[
  \phi(x)-\phi(y)-\phi'(y)(x-y)=x\log(x/y)+y-x.
\]
Convexity makes it nonnegative for $x\ge0$. On probability laws the $y-x$
terms sum to zero, leaving the relative entropy
\[
  D(\bar p_i\|\bar q_{-i})
    =\sum_h\bar p_i(h)\log\frac{\bar p_i(h)}{\bar q_{-i}(h)}.
\]
Each projected convex term has a nonnegative remainder, and the linear part
of the full objective expansion vanishes because the gradient at
$\bar q_{-i}$ is constant and the two laws have the same mass. Therefore
$\tau D(\bar p_i\|\bar q_{-i})\le\mathcal F_i(\bar p_i)-\mathcal F_i(\bar q_{-i})$.
For probability laws $p,r$ with positive reference law $r$,
$-\log x\ge1-x$ at $x=\sqrt{r(h)/p(h)}$ gives
$p(h)\log(p(h)/r(h))\ge2p(h)-2\sqrt{p(h)r(h)}$ for $p(h)>0$.
Zero $p(h)$ contributes zero on both sides. Summing proves $H^2(p,r)\le D(p\|r)$.
Here $\bar q_{-i}$ is positive on its realized alphabet by the positivity
argument above. Since $e_i$ is injective,
\begin{equation}
  H^2((e_i)_\#\bar p_i,q_{-i})\le(\log2)\mu_i\le\mu_i.
  \label{eq:deletion-kl-cost}
\end{equation}

On a pair with total mass $u$, replacing the equal split by all mass on the
actual representative contributes
$(\sqrt{u/2}-\sqrt{u})^2+u/2=(2-\sqrt 2)u$ to $H^2$. Singletons are
unchanged, so $H^2(r_i',(e_i)_\#\bar p_i)=(2-\sqrt2)\mu_i\le\mu_i$.
The inequality $H^2(p,r)\le2\bigl(H^2(p,v)+H^2(v,r)\bigr)$ follows by
applying the Euclidean inequality $\|x+y\|^2\le2(\|x\|^2+\|y\|^2)$ to
square-root vectors. With \eqref{eq:deletion-kl-cost} it yields
\begin{equation}
  H^2(r_i',q_{-i})\le 4\mu_i.
  \label{eq:deletion-collapse-cost}
\end{equation}

For every two-element fiber mark its smaller-mass endpoint under $q$, with
an arbitrary tie convention. Let $E_i$ be the marked set. Then
$r_i'(E_i)=\mu_i/2$, while $\sum_i q(E_i)\le d$ by
\eqref{eq:deletion-input-edges}, since the two-element fibers at $i$ are
exactly the $i$-edges. For any $E\subseteq G$ the scalar inequality
$(x+y)^2\le2x^2+2y^2$ at $x=\sqrt{p(h)}$ and
$y=\sqrt{r(h)}-\sqrt{p(h)}$ gives $r(E)\le2p(E)+2H^2(p,r)$.
It follows that
\begin{equation}
  \sum_i\mu_i\le4d+4\sum_i H^2(q,r_i').
  \label{eq:deletion-marked-mass}
\end{equation}

\subsection{The summed deletion comparison}

\begin{proposition}[Deletion on a common block]
\label{prop:deletion-summed}
For fixed weights, let $q$, $r_i$, and $q_{-i}$ have the meanings above. Then
\begin{equation}
  \sum_i H^2(q,q_{-i})\le576d+68\sum_i H^2(q,r_i).
  \label{eq:deletion-pointwise-budget}
\end{equation}
For independent uniform weights on $[0,2]$, its expectation is at most
$4928d$. The same bound holds for the deletion distances between the actual
output laws of the randomized learner, after averaging the weights.
\end{proposition}

\begin{proof}
Set $S=\sum_i H^2(q,r_i')$ and $L=\sum_i H^2(q,r_i)$. The squared triangle
inequality and \eqref{eq:deletion-restriction-cost} give $S\le2(L+NB)$.
Applying the same inequality through $r_i'$ and substituting
\eqref{eq:deletion-collapse-cost} and \eqref{eq:deletion-marked-mass} gives
\begin{align*}
  \sum_i H^2(q,q_{-i})
     &\le 2S+8\sum_i\mu_i\le32d+34S \\
     &\le32d+68(L+NB)\le576d+68L,
\end{align*}
using $NB\le8d$ from \eqref{eq:deletion-restriction-budget}. Taking
expectations and applying \eqref{eq:deletion-loss-budget} bounds the result
by $(576+68\cdot64)d=4928d$.

To compare the actual output laws, mix $q$ and $q_{-i}$ using the same full
block weight vector. The unused weight $U_i$ has no effect on $q_{-i}$, so
both marginals are the intended output laws. For any coupled random laws
$p_U,r_U$, Cauchy--Schwarz gives
$\mathbb E_U\sqrt{p_U(h)r_U(h)}\le\sqrt{\mathbb E_U p_U(h)\,\mathbb E_U r_U(h)}$
for each $h$. Expanding the squared Hellinger distance and summing over $h$
yields
$H^2(\mathbb E_U p_U,\mathbb E_U r_U)\le\mathbb E_U H^2(p_U,r_U)$.
This proves the last claim. The bound $4928d$ can now be used for the
common-block deletion sum in the information argument.
\end{proof}


\section{Correspondence with the formal proof}\label{app:map}
The Lean kernel checks the formal propositions under their stated hypotheses. The table
below gives the correspondence needed to review their interpretation as a
learning experiment. Section~\ref{sec:scope} compares the resulting theorem
with the original conjecture. Both Borel obstructions in that section are
formalized, including the existence of a continuous map with non-Borel range.

The declarations below belong to the namespace \texttt{AgnosticVC}.
The \href{https://github.com/1337isnot1337/solved-conjectures/tree/fc810afee89754366e8e13c8bb8a847c624c3a8c/lean/Conjectures/AgnosticVC}{pinned source of the rate proof}
and the \href{https://github.com/1337isnot1337/solved-conjectures/actions/runs/44981182433}{public replay of the same Lean sources}
document the original 61-module development. An \href{https://github.com/1337isnot1337/solved-conjectures/actions/runs/44979669007}{exact-source run}
also verifies the linked proof commit. The run elaborated the sources
with warnings treated as errors, audited 1,106 declarations, and replayed
the compiled proofs. The only transitive axioms in that audit were
\texttt{propext}, \texttt{Classical.choice}, and \texttt{Quot.sound}.
The subsequent results on measurable selection and its obstructions are in the
\href{https://github.com/1337isnot1337/solved-conjectures/tree/35970639ec69a753950f4ed663f9a8c704f4fca4/lean/Conjectures/AgnosticVC}{extended development},
with a \href{https://github.com/1337isnot1337/solved-conjectures/actions/runs/44986814421}{separate verification run}.
The obstruction applies to arbitrary proper Markov kernels, including kernels
whose output laws have infinite support.

{\small\setlength{\tabcolsep}{4pt}\renewcommand{\arraystretch}{1.13}
\begin{longtable}{@{}>{\raggedright\arraybackslash}p{.36\linewidth}>{\raggedright\arraybackslash}p{.60\linewidth}@{}}
\toprule Mathematical step & Lean declaration(s) \\ \midrule \endhead
Global names survive restriction & \decl{canonicalRep_extendTrace} \\
Optimization and sensitivity & \decl{minimizer_unique}, \decl{optimizer_pos}, \decl{optimizer_hasDerivWithinAt} \\
Trace entropy and packing & \decl{log_card_traceRestriction_le}, \decl{log_card_le_of_hamming_separated} \\
Random weights and all-sample accuracy & \decl{signAverage_max_binaryLoss_le}, \decl{integral_maxWeightDeviation_le}, \decl{properLearner_empirical_le_iInf} \\
Projection overlap and inverse & \decl{averagedProjection_sqrt_overlap}, \decl{projectionScale_schur}, \decl{projectionScale_inverse_bound} \\
Weighted edges and residual variance & \decl{weighted_oneInclusionEdges_le}, \decl{traceProjection_average_residual_le} \\
Score and loss removal & \decl{traceProjection_hessian}, \decl{traceProjection_scoreEnergy_sum}, \decl{traceLaw_remove_loss_integral_le} \\
Restriction replacement and collapse & \decl{optimizer_remove_restrictions}, \decl{optimizer_collapse_relativeEntropy}, \decl{sum_doubleMass_forgetTrace_le} \\
Total deletion cost & \decl{sum_hellingerSq_schedule_deleted_le}, \decl{blockDeletion_sum_le} \\
Whole-output KL and bit comparison & \decl{poolInformation_eq_klDiv}, \decl{poolInformation_le_blockCharges} \\
Fixed-pool and iid exchangeability & \decl{steinkeZakynthinou_fixedPool}, \decl{steinkeZakynthinou_measurableInformation} \\
Measurable Markov learner & \decl{measurable_properLearner_toMeasure}, \decl{steinkeZakynthinou_conjecture6} \\
Countable instance domain & \decl{steinkeZakynthinou_countable_domain} \\
Countable initial segment realizing all traces & \decl{hasMeasurableCanonicalRepresentatives_of_countable_trace_core} \\
Ordered empirical risk minimization & \decl{hasMeasurableCanonicalRepresentatives_of_orderedERM}, \decl{steinkeZakynthinou_of_measurable_orderedERM} \\
Choice of a suitable well-order & \decl{exists_order_of_countable_trace_core} \\
Pointwise limits and continuous parameters & \decl{finite_trace_core_of_pointwise_limits}, \decl{exists_order_of_separable_continuous_predict} \\
Rational cuts for real thresholds & \decl{rational_threshold_realizes_finite_trace} \\
Repeated positive observations & \decl{measurableSet_exists_predict_true_of_kernel_rate} \\
Singleton classifiers on the unit interval & \decl{exists_unitInterval_singleton_obstruction} \\
Continuous map with non-Borel range & \decl{exists_continuous_nat_nat_nonmeasurable_range} \\
Standard Borel counterexample & \decl{exists_standardBorel_vcOne_no_accurateKernel}, \decl{borelPair_no_accurateKernel_rate} \\
\bottomrule
\end{longtable}}

\section*{Acknowledgments}
The project used AI assistance in preparing the manuscript and Lean development.

{\small
\begin{thebibliography}{9}
\setlength{\itemsep}{1pt}
\bibitem{sauer} N. Sauer.
\href{https://doi.org/10.1016/0097-3165(72)90019-2}{On the density of families of sets}.
\emph{Journal of Combinatorial Theory, Series A} 13 (1972), 145--147.
\bibitem{hlw} D. Haussler, N. Littlestone and M. K. Warmuth.
\href{https://mwarmuth.bitbucket.io/pubs/C10.pdf}{Predicting $\{0,1\}$-functions on randomly drawn points}.
29th Annual Symposium on Foundations of Computer Science, 1988, 100--109.
\bibitem{problem} T. Steinke and L. Zakynthinou.
\href{https://proceedings.mlr.press/v125/steinke20b.html}{Open Problem: Information Complexity of VC Learning}.
Conference on Learning Theory, 2020.
\bibitem{framework} T. Steinke and L. Zakynthinou.
\href{https://proceedings.mlr.press/v125/steinke20a.html}{Reasoning About Generalization via Conditional Mutual Information}.
Conference on Learning Theory, 2020. Definition and lemma numbers refer to the
\href{https://arxiv.org/abs/2001.09122}{full version}.
\bibitem{haghifam} M. Haghifam, G. K. Dziugaite, S. Moran and D. M. Roy.
\href{https://arxiv.org/abs/2011.02970}{On the Information Complexity of Proper Learners for VC Classes in the Realizable Case}, 2020.
\bibitem{hanneke} Steve Hanneke and Juexiao Wang.
\href{https://arxiv.org/abs/2610.10600}{The Optimal Information Complexity of VC Learning}, 2026.
\bibitem{kechris} A. S. Kechris.
\href{https://doi.org/10.1007/978-1-4612-4190-4}{Classical Descriptive Set Theory}.
Graduate Texts in Mathematics 156, Springer, 1995.
\end{thebibliography}}
\end{document}
